\subsection{\texorpdfstring{L.4 The $k\to\ell$ Gate (CMB Observability as Policy Interface)}{L.4 The k->ell Gate (CMB Observability as Policy Interface)}}
\begin{tcolorbox}[colback=yellow!3, colframe=yellow!60!black,
    title={\textbf{What is deterministic vs.\ what is still policy}}, fonttitle=\bfseries\small]
\textbf{Deterministic (given the bounce completion):}
\begin{itemize}[leftmargin=1.5em, topsep=2pt, itemsep=1pt]
    \item The transfer-function shape $T(k)$ (this appendix; implemented in \texttt{bounce\_perturbations}).
    \item Given an explicit expansion-history budget, the mapping $k\to\ell$ and the resulting $C_\ell$ curves are deterministic (implemented in \texttt{k\_calibration} and \texttt{boltzmann\_transfer}).
\end{itemize}

\textbf{Policy / not yet ToE-closed:}
\begin{itemize}[leftmargin=1.5em, topsep=2pt, itemsep=1pt]
    \item The mapping from comoving $k$ to observed CMB multipole $\ell$ depends on the absolute expansion history (reheating/threshold bookkeeping), summarized as an overall scaling $a_0/a_t$ (or more precisely $a_0/a_{\rm transition}$ for the chosen transition surface).
    \item The signature classification (``CMB'' vs ``small-scale'') is therefore a \emph{gate}: it depends on the declared history, and should not be framed as an unconditional claim.
\end{itemize}

\textbf{Notation contract (to prevent $k$-confusion):}
\begin{itemize}[leftmargin=1.5em, topsep=2pt, itemsep=1pt]
    \item $k_{\hat{}}$ is a \emph{dimensionless} wave-number used inside the bounce solver (in $x=M\eta$ units).
    \item The raw bounce scales $k^{(s)}_{\rm bounce}$ and $k^{(t)}_{\rm bounce}$ are reported in these internal units; absolute $k(\mathrm{Mpc}^{-1})$ requires the TFPT scale $M$ and the expansion budget $a_0/a_{\rm transition}$.
    \item The observable proxy is $\ell \approx k(\mathrm{Mpc}^{-1})\,\chi_\*$ with $\chi_\*$ computed in a flat-$\Lambda$CDM snapshot (suite).
\end{itemize}

\textbf{Current suite snapshot (physics mode; threshold-derived budget):}
\begin{itemize}[leftmargin=1.5em, topsep=2pt, itemsep=1pt]
    \item $a_0/a_{\rm transition}\approx 2.4486\times 10^{57}$,
    \item $\chi_\*\approx 1.3867\times 10^4$ Mpc,
    \item $\ell^{(s)}_{\rm bounce}\approx 56.997$ and $\ell^{(t)}_{\rm bounce}\approx 0.212$.
\end{itemize}
Under this explicit gate, the scalar bounce feature is CMB-low-$\ell$, while the tensor feature is super-horizon (below the quadrupole). 

\textbf{Falsifiable signature classes (conditional, gate-dependent):}
\begin{itemize}[leftmargin=1.5em, topsep=2pt, itemsep=1pt]
    \item \textbf{CMB-low-$\ell$ scalar feature}: features near $\ell\sim\mathcal{O}(10^2)$ (current default gate).
    \item \textbf{Small-scale-only}: features shifted to $\ell\gg 10^3$ under alternative histories.
\end{itemize}

\textbf{Suite evidence:} The \texttt{k\_calibration} module reports both (i) the required scaling to place features at target multipoles and (ii) the threshold-derived expansion-budget estimate; \texttt{boltzmann\_transfer} produces explicit CAMB-backed $C_\ell$ curves and the relative deviation $\Delta C_\ell/C_\ell$ between bounce-injected and baseline spectra.
\end{tcolorbox}

\TFPTkCalibrationTable

\begin{figure}[H]
\centering
\begin{tikzpicture}
\begin{groupplot}[
    group style={group size=2 by 1, horizontal sep=1.2cm},
    width=0.48\textwidth,
    height=0.42\textwidth,
    grid=both,
    legend style={font=\scriptsize},
]
\nextgroupplot[
    title={required $a_0/a_{\rm tr}$ vs target $\ell$},
    xmode=log,
    ymode=log,
    xlabel={target $\ell$},
    ylabel={$a_0/a_{\rm tr}$},
    xmin=1.5, xmax=1000,
]
\addplot+[mark=*] coordinates {
    (2,6.978134315108868e58)
    (30,4.652089543405912e57)
    (700,1.9937526614596766e56)
};
\addlegendentry{scalar (required)}
\addplot+[mark=square*] coordinates {
    (2,2.5901704412653403e56)
    (30,1.7267802941768936e55)
    (700,7.400486975043829e53)
};
\addlegendentry{tensor (required)}
\addplot[red, dashed, thick] coordinates {(2,2.4485807685210785e57) (700,2.4485807685210785e57)};
\addlegendentry{current budget ($2.45\times10^{57}$)}

\nextgroupplot[
    title={gate view (CMB window + bounce $\ell$)},
    xmode=log,
    xlabel={$\ell$},
    ymin=0, ymax=1,
    ytick=\empty,
    xmin=0.1, xmax=5000,
    axis y line=none,
]
% CMB window shading: [2,2500]
\addplot[draw=none, fill=blue!10] coordinates {(2,0) (2500,0) (2500,1) (2,1)} \closedcycle;
\addplot[black!40] coordinates {(2,0) (2,1)};
\addplot[black!40] coordinates {(2500,0) (2500,1)};
\node[anchor=south, font=\scriptsize] at (axis cs:20,0.02) {CMB window};

% bounce locations (physics-suite snapshot)
\addplot[green!60!black, thick] coordinates {(56.99737909257206,0) (56.99737909257206,1)};
\node[anchor=south, font=\scriptsize, green!60!black] at (axis cs:56.99737909257206,0.02) {$\ell_b^{(s)}\approx57$};
\addplot[orange!80!black, thick] coordinates {(0.2115650400072186,0) (0.2115650400072186,1)};
\node[anchor=south, font=\scriptsize, orange!80!black] at (axis cs:0.2115650400072186,0.02) {$\ell_b^{(t)}\approx0.21$};
\end{groupplot}
\end{tikzpicture}
\caption{Suite plots from \texttt{k\_calibration} (physics mode): scaling and feasibility diagnostics for the $k\to\ell$ gate under an explicit expansion-history budget.}
\end{figure}

\begin{figure}[H]
\centering
\begin{tikzpicture}
\begin{groupplot}[
    group style={group size=2 by 2, horizontal sep=1.2cm, vertical sep=1.0cm},
    width=0.48\textwidth,
    height=0.33\textwidth,
    xmode=log,
    xmin=2, xmax=3000,
    grid=both,
    legend style={font=\scriptsize},
]
\nextgroupplot[
    title={$D_\ell^{TT}$ (baseline vs bounce)},
    ymode=log,
    ylabel={$D_\ell^{TT}$ [$\mu$K$^2$]},
]
\addplot+[mark=*, mark size=2pt] coordinates {
    (2,1053.602432)
    (3,998.0887611)
    (4,944.9109792)
    (5,905.3057583)
    (10,845.7614033)
    (20,936.246946)
    (30,1090.649646)
    (50,1468.079395)
    (80,2185.897476)
    (100,2782.530282)
    (150,4517.548558)
    (200,5768.991349)
    (300,4189.951412)
    (500,2521.115802)
    (700,1928.629066)
    (1000,1093.569344)
    (1500,706.1670601)
    (2000,237.9250741)
    (2508,78.09810466)
};
\addlegendentry{baseline}
\addplot+[mark=square*, mark size=2pt] coordinates {
    (2,10994.05279)
    (3,10486.33276)
    (4,9385.765294)
    (5,7182.088179)
    (10,466.5297095)
    (20,311.2303573)
    (30,647.8609341)
    (50,1396.409529)
    (80,2178.001377)
    (100,2764.125993)
    (150,4513.504899)
    (200,5742.64759)
    (300,4165.032244)
    (500,2519.57907)
    (700,1945.981149)
    (1000,1115.817699)
    (1500,728.4072936)
    (2000,247.883187)
    (2508,81.98087433)
};
\addlegendentry{bounce injected}

\nextgroupplot[
    title={$D_\ell^{EE}$ (baseline vs bounce)},
    ymode=log,
    ylabel={$D_\ell^{EE}$ [$\mu$K$^2$]},
]
\addplot+[mark=*, mark size=2pt] coordinates {
    (2,0.03176782281)
    (3,0.0407463901)
    (4,0.03535540342)
    (5,0.02363620444)
    (10,0.003157610827)
    (20,0.005394191594)
    (30,0.0223578033)
    (50,0.1235977312)
    (80,0.4863197226)
    (100,0.8033626171)
    (150,1.135920327)
    (200,0.710593948)
    (300,9.377843402)
    (500,8.525653848)
    (700,38.4094504)
    (1000,43.34908241)
    (1500,13.22941615)
    (2000,8.84123834)
    (2508,3.009426377)
};
\addlegendentry{baseline}
\addplot+[mark=square*, mark size=2pt] coordinates {
    (2,0.5094182832)
    (3,0.562182037)
    (4,0.3663058727)
    (5,0.1665671877)
    (10,0.0005850428195)
    (20,0.0006453541167)
    (30,0.004809524067)
    (50,0.09804048815)
    (80,0.489652682)
    (100,0.7951178048)
    (150,1.13708347)
    (200,0.7069214332)
    (300,9.311700761)
    (500,8.492435159)
    (700,38.50920372)
    (1000,43.9978482)
    (1500,13.62199765)
    (2000,9.189186049)
    (2508,3.154751105)
};
\addlegendentry{bounce injected}

\nextgroupplot[
    title={relative deviation (TT)},
    xlabel={$\ell$},
    ylabel={$\Delta D_\ell/D_{\ell,0}$},
    ymin=-1.0, ymax=10.5,
]
\addplot+[mark=*, mark size=2pt] coordinates {
    (2,9.43472609)
    (3,9.506413022)
    (4,8.93296247)
    (5,6.933328727)
    (10,-0.4483908727)
    (20,-0.6675766381)
    (30,-0.405986206)
    (50,-0.0488187938)
    (80,-0.003612291319)
    (100,-0.006614227665)
    (150,-0.0008951004397)
    (200,-0.004566441139)
    (300,-0.005947364325)
    (500,-0.0006095444869)
    (700,0.008997107253)
    (1000,0.02034471301)
    (1500,0.03149429467)
    (2000,0.04185398693)
    (2508,0.04971656719)
};
\addplot[black!50, dashed] coordinates {(2,0) (2508,0)};

\nextgroupplot[
    title={relative deviation (EE)},
    xlabel={$\ell$},
    ylabel={$\Delta D_\ell/D_{\ell,0}$},
    ymin=-1.0, ymax=16.5,
]
\addplot+[mark=*, mark size=2pt] coordinates {
    (2,15.03566874)
    (3,12.79710045)
    (4,9.360675802)
    (5,6.04712079)
    (10,-0.8147197829)
    (20,-0.8803612913)
    (30,-0.7848838724)
    (50,-0.2067776069)
    (80,0.006853432569)
    (100,-0.01026287768)
    (150,0.00102396489)
    (200,-0.005168232629)
    (300,-0.007053075852)
    (500,-0.003896321651)
    (700,0.002597103541)
    (1000,0.0149660788)
    (1500,0.02967489204)
    (2000,0.0393550876)
    (2508,0.04828984317)
};
\addplot[black!50, dashed] coordinates {(2,0) (2508,0)};
\end{groupplot}
\end{tikzpicture}
\caption{Suite plots from \texttt{boltzmann\_transfer} (physics mode): CAMB-backed $C_\ell$ curves (TT/EE) for baseline vs bounce-injected primordial spectrum, and the relative deviation $\Delta C_\ell/C_\ell$.}
\end{figure}

\subsection{L.5 Falsification Targets}
Large-scale suppression or oscillatory features in $\mathcal{P}_\zeta(k)$ are predicted at
$k\lesssim k_{\text{bounce}}$. Current CMB large-scale bounds constrain the allowable feature
amplitude, providing a direct observational test of the bounce completion.

%=============================================================================
\section*{Reference Ledger (Auto-Generated)}
%=============================================================================
\addcontentsline{toc}{section}{Reference Ledger (Auto-Generated)}
% Inlined (standalone build): formerly \input{../tfpt-suite/out/paper_reference_ledger.tex}
% -----------------------------------------------------------------------------
% AUTO-GENERATED by tfpt_suite/export_reference_ledger.py — do not edit by hand.
% generated_at_utc: 2026-01-28T07:36:57.179559+00:00
% source: tfpt_suite/data/references.json (+ dereferenced global_reference*.json observables where applicable)

\begingroup
\footnotesize
\setlength{\tabcolsep}{4pt}
\renewcommand{\_}{\string_\allowbreak}
\Urlmuskip=0mu plus 1mu\relax
\begin{longtable}{@{}p{0.24\textwidth}rrllp{0.34\textwidth}@{}}
\caption{Reference ledger used by the TFPT suite (auto-generated).}\label{tab:reference-ledger}\\
\toprule
Key & Value & $\sigma$ & Units & Version & Source (stable key)\\
\midrule
\endfirsthead
\toprule
Key & Value & $\sigma$ & Units & Version & Source (stable key)\\
\midrule
\endhead
\midrule
\multicolumn{6}{r}{\small Continued on next page.}\\
\midrule
\endfoot
\bottomrule
\endlastfoot
\texttt{A\_s\_planck2018} & \num{2.0989031673e-09} & \num{2.938464434e-11} & dimensionless & Planck 2018 & \cite{planck2018} Planck 2018 results VI (2020), base-$\Lambda$CDM TT,TE,EE+lowE+lensing: ln(10\textasciicircum{}\{10\} A\_s)=3.044 $\pm$ 0.014 $\Rightarrow$ A\_s=exp(3.044)/1e10 (\url{https://arxiv.org/abs/1807.06209})\\
\texttt{H0\_planck2018} & \num{67.36} & \num{0.54} & \(\mathrm{km\,s^{-1}\,Mpc^{-1}}\) & Planck 2018 & \cite{planck2018} Planck 2018 results VI (2020), base-$\Lambda$CDM TT,TE,EE+lowE+lensing: $H_0 = 67.36 \pm 0.54$ (Table 2) (\url{https://arxiv.org/abs/1807.06209})\\
\texttt{H0\_sh0es\_2022} & \num{73.04} & \num{1.04} & \(\mathrm{km\,s^{-1}\,Mpc^{-1}}\) & SH0ES 2022 & \cite{riess2022} Riess et al. 2022 (ApJ 934 L7)\\
\texttt{Mpl\_reduced} & \num{2.435e+18} & \num{0.0} & \(\mathrm{GeV}\) & CODATA & \cite{tfptSuiteData} tfpt\_suite.cosmo\_scale\_map.MPL\_REDUCED\_GEV\\
\texttt{Omega\_dm\_h2\_planck2018} & \num{0.12} & \num{0.001} & dimensionless & Planck 2018 & \cite{planck2018} Planck 2018 base LCDM\\
\texttt{alpha\_bar5\_inv\_MZ\_pdg2024} & \num{127.93} & \num{0.008} & dimensionless & PDG 2024 & \cite{pdg2024} PDG 2024 electroweak review: $\hat{\alpha}^{(5)}(M_Z)^{-1}$ = 127.930 $\pm$ 0.008 (MSbar; see tfpt\_suite/data/alpha\_running\_pdg.json).\\
\texttt{alpha\_inv\_codata\_2022} & \num{137.035999177} & \num{2.1e-08} & dimensionless & CODATA 2022 & \cite{codata2022} NIST CODATA 2022: inverse fine-structure constant $\alpha^{-1}$=137.035999177(21) (\url{https://physics.nist.gov/cgi-bin/cuu/Value?alphinv})\\
\texttt{alpha\_s\_mz\_pdg} & \num{0.1179} & \num{0.0011} & dimensionless & PDG (MZ, MSbar) & \cite{pdg2024} sm\_inputs\_mz.json:alpha\_s\\
\texttt{beta\_deg\_minami\_komatsu\_2020} & \num{0.35} & \num{0.14} & \(\mathrm{deg}\) & Minami \& Komatsu 2020 & \cite{minamiKomatsu2020} Minami \& Komatsu (2020), Phys. Rev. Lett. 125, 221301: $\beta = 0.35^\circ \pm 0.14^\circ$ (Planck 2018 polarization, accounting for calibration) (\url{https://doi.org/10.1103/PhysRevLett.125.221301})\\
\texttt{cabibbo\_lambda\_pdg2024} & \num{0.22501} & \num{0.00068} & dimensionless & PDG 2024 & \cite{pdg2024} PDG 2024 CKM review (global fit, Eq. 12.27): $|V_{us}| = 0.22501 \pm 0.00068$ (\url{https://pdg.lbl.gov/2024/reviews/rpp2024-rev-ckm-matrix.pdf})\\
\texttt{ln10\_As\_planck2018} & \num{3.044} & \num{0.014} & dimensionless & Planck 2018 & \cite{planck2018} Planck 2018 results VI (2020), base-$\Lambda$CDM TT,TE,EE+lowE+lensing: ln(10\textasciicircum{}\{10\} A\_s) = 3.044 $\pm$ 0.014 (Table 2) (\url{https://arxiv.org/abs/1807.06209})\\
\texttt{m\_b\_pdg} & \num{4.18} & --- & \(\mathrm{GeV}\) & PDG (scheme dependent) & \cite{pdg2024} sm\_inputs\_mz.json:mb\_GeV\\
\texttt{m\_c\_pdg} & \num{1.27} & --- & \(\mathrm{GeV}\) & PDG (scheme dependent) & \cite{pdg2024} sm\_inputs\_mz.json:mc\_GeV\\
\texttt{m\_e\_pdg} & \num{0.0005109989461} & \num{0.0} & \(\mathrm{GeV}\) & PDG pole mass & \cite{pdg2024} lepton\_masses\_pdg.json:masses.electron\\
\texttt{m\_mu\_pdg} & \num{0.1056583745} & \num{0.0} & \(\mathrm{GeV}\) & PDG pole mass & \cite{pdg2024} lepton\_masses\_pdg.json:masses.muon\\
\texttt{m\_p\_pdg} & \num{0.938272} & --- & \(\mathrm{GeV}\) & PDG pole mass & \cite{pdg2024} mass\_spectrum\_minimal:ledger.placeholders.proton\_mass\_GeV\\
\texttt{m\_t\_pdg} & \num{172.76} & \num{0.3} & \(\mathrm{GeV}\) & PDG pole mass & \cite{pdg2024} sm\_inputs\_mz.json:mt\_GeV\\
\texttt{m\_tau\_pdg} & \num{1.77686} & \num{0.0} & \(\mathrm{GeV}\) & PDG pole mass & \cite{pdg2024} lepton\_masses\_pdg.json:masses.tau\\
\texttt{mass\_ratio\_mb\_over\_ms} & --- & --- & dimensionless & PDG (scheme dependent) & \cite{tfptSuiteData} \\
\texttt{mass\_ratio\_mc\_over\_mu\_quark} & --- & --- & dimensionless & PDG (scheme dependent) & \cite{tfptSuiteData} \\
\texttt{mass\_ratio\_ms\_over\_md} & --- & --- & dimensionless & PDG (scheme dependent) & \cite{tfptSuiteData} \\
\texttt{mass\_ratio\_mt\_over\_mc} & --- & --- & dimensionless & PDG (scheme dependent) & \cite{tfptSuiteData} \\
\texttt{mass\_ratio\_mu\_over\_e\_pdg} & \num{206.76828260879265} & \num{0.0} & dimensionless & PDG pole masses & \cite{pdg2024} lepton\_masses\_pdg.json (mu/e)\\
\texttt{mass\_ratio\_tau\_over\_mu\_pdg} & \num{16.81702949159037} & \num{0.0} & dimensionless & PDG pole masses & \cite{pdg2024} lepton\_masses\_pdg.json (tau/mu)\\
\texttt{n\_s\_planck2018} & \num{0.9649} & \num{0.0042} & dimensionless & Planck 2018 & \cite{planck2018} Planck 2018 results VI (2020), base-$\Lambda$CDM TT,TE,EE+lowE+lensing: $n_s = 0.9649 \pm 0.0042$ (Table 2) (\url{https://arxiv.org/abs/1807.06209})\\
\texttt{omega\_b\_h2\_planck2018} & \num{0.02237} & \num{0.00015} & dimensionless & Planck 2018 & \cite{planck2018} Planck 2018 results VI (2020), base-$\Lambda$CDM TT,TE,EE+lowE+lensing: $\Omega_b h^2 = 0.02237 \pm 0.00015$ (Table 2) (\url{https://arxiv.org/abs/1807.06209})\\
\texttt{pmns\_delta\_cp\_deg\_nufit53\_no} & \num{232.0} & \num{39.0} & \(\mathrm{deg}\) & NuFIT 5.3 (2024, NO) & \cite{nufit53_2024} pmns\_reference.json:normal\_ordering.delta\_cp\_deg\\
\texttt{pmns\_sin2\_theta12\_nufit53\_no} & \num{0.307} & \num{0.012} & dimensionless & NuFIT 5.3 (2024, NO) & \cite{nufit53_2024} pmns\_reference.json:normal\_ordering.sin2\_theta12\\
\texttt{pmns\_sin2\_theta13\_nufit53\_no} & \num{0.02224} & \num{0.00057} & dimensionless & NuFIT 5.3 (2024, NO) & \cite{nufit53_2024} pmns\_reference.json:normal\_ordering.sin2\_theta13\\
\texttt{pmns\_sin2\_theta23\_nufit53\_no} & \num{0.454} & \num{0.019} & dimensionless & NuFIT 5.3 (2024, NO) & \cite{nufit53_2024} pmns\_reference.json:normal\_ordering.sin2\_theta23\\
\texttt{sin2\_thetaW\_mz\_pdg} & \num{0.23122} & \num{4e-05} & dimensionless & PDG (MZ, MSbar) & \cite{pdg2024} sm\_inputs\_mz.json:sin2\_thetaW\\
\texttt{v\_ew\_246} & \num{246.0} & \num{0.2} & \(\mathrm{GeV}\) & SM Higgs vev & \cite{tfptSuiteData} Standard Model reference\\
\end{longtable}
\endgroup

%=============================================================================
\section*{References}
%=============================================================================

\begin{thebibliography}{99}
\bibitem{hehl1976}
F.~W.~Hehl, P.~von der Heyde, G.~D.~Kerlick, J.~M.~Nester,
``General Relativity with Spin and Torsion: Foundations and Prospects,''
Rev.\ Mod.\ Phys.\ \textbf{48}, 393 (1976).

\bibitem{fujikawa1979}
K.~Fujikawa,
``Path-Integral Measure for Gauge-Invariant Fermion Theories,''
Phys.\ Rev.\ Lett.\ \textbf{42}, 1195 (1979).

\bibitem{kawasaki1979}
T.~Kawasaki,
``The Riemann--Roch theorem for complex V-manifolds,''
Osaka J.\ Math.\ \textbf{16}, 151--159 (1979).

\bibitem{atiyah1975}
M.~F.~Atiyah, V.~K.~Patodi, I.~M.~Singer,
``Spectral asymmetry and Riemannian geometry I,''
Math.\ Proc.\ Camb.\ Phil.\ Soc.\ \textbf{77}, 43--69 (1975).

\bibitem{bunke1995}
U.~Bunke,
``On the gluing problem for the $\eta$-invariant,''
J.\ Differential Geom.\ \textbf{41}, 397--448 (1995).

\bibitem{diegoPalazuelos2022}
P.~Diego-Palazuelos \textit{et al.},
``Cosmic Birefringence from the Planck Data Release 4,''
Phys.\ Rev.\ Lett.\ \textbf{128}, 091302 (2022).

\bibitem{minamiKomatsu2020}
Y.~Minami, E.~Komatsu,
``New Extraction of the Cosmic Birefringence from the Planck 2018 Polarization Data,''
Phys.\ Rev.\ Lett.\ \textbf{125}, 221301 (2020).

\bibitem{codata2022}
E.~Tiesinga, P.~J.~Mohr, D.~B.~Newell, B.~N.~Taylor,
``CODATA Recommended Values of the Fundamental Physical Constants: 2022,''
J.\ Phys.\ Chem.\ Ref.\ Data \textbf{53}, 043103 (2024).

\bibitem{planck2018}
Planck Collaboration,
``Planck 2018 results. VI. Cosmological parameters,''
\textit{Astron.\ Astrophys.} \textbf{641}, A6 (2020), arXiv:1807.06209.

\bibitem{pdg2024}
Particle Data Group,
``Review of Particle Physics (2024),''
see PDG 2024 online review and summary tables.

\bibitem{nufit53_2024}
NuFIT collaboration,
``NuFIT 5.3 (2024) global analysis of neutrino oscillation data,''
see \texttt{http://www.nu-fit.org/} (snapshot 2024).

\bibitem{riess2022}
A.~G.~Riess \textit{et al.},
``A Comprehensive Measurement of the Local Value of the Hubble Constant with 1 km s$^{-1}$ Mpc$^{-1}$ Uncertainty from the Hubble Space Telescope and the SH0ES Team,''
Astrophys.\ J.\ Lett.\ \textbf{934}, L7 (2022).

\bibitem{kostelecky2008}
V.~A.~Kosteleck\'y, N.~Russell, J.~D.~Tasson,
``Constraints on Torsion from Bounds on Lorentz Violation,''
Phys.\ Rev.\ Lett.\ \textbf{100}, 111102 (2008).

\bibitem{shapiro2002}
I.~L.~Shapiro,
``Physical aspects of the space--time torsion,''
Phys.\ Rept.\ \textbf{357}, 113--213 (2002).

\bibitem{tfptSuiteData}
TFPT Suite (v2.5) data + policy files,
\codepath{tfpt-suite/tfpt_suite/data/} and \codepath{tfpt-suite/out/} artifacts (this repository; 2026-01-27 snapshot).
\end{thebibliography}

\end{document}
